%% Cover Letter — generated by generate-letter.mjs

\documentclass[a4paper, 10.5pt]{article}
\usepackage[top=28mm, bottom=24mm, left=28mm, right=28mm]{geometry}
\usepackage{xcolor}
\usepackage{tikz}
\usepackage[hidelinks]{hyperref}
\usepackage[french]{babel}
\usepackage[T1]{fontenc}
\usepackage[sfdefault]{inter}
\usepackage{raleway}
\usepackage{parskip}

\pagestyle{empty}
\setlength{\parindent}{0pt}
\setlength{\parskip}{7pt}

\tolerance=200
\emergencystretch=2em

\newcommand{\headingfont}{\fontfamily{Raleway-TLF}\selectfont}

% Colors
\definecolor{accent}{HTML}{B5651D}
\definecolor{titletext}{HTML}{111827}
\definecolor{bodytext}{HTML}{374151}
\definecolor{mutedtext}{HTML}{6B7280}
\definecolor{sep}{HTML}{D1D5DB}
\definecolor{rulecolor}{HTML}{E5E7EB}

\begin{document}
\color{bodytext}

% ─── Header ──────────────────────────────────────────────────────────────────
\begin{center}
{\headingfont\fontsize{20}{24}\selectfont\bfseries\color{titletext}Amaury Dufrenot}\\[4pt]
{\small\color{mutedtext}Étudiant ingénieur Réseaux \& Télécommunications — UTT}\\[6pt]
{\small\color{mutedtext}\href{mailto:amaury.dufrenot@utt.fr}{amaury.dufrenot@utt.fr}\hspace{6pt}{\color{sep}\textbar}\hspace{6pt}+33 6 44 19 68 50\hspace{6pt}{\color{sep}\textbar}\hspace{6pt}\href{https://amaury-dufrenot.me}{amaury-dufrenot.me}}\\[8pt]
{\color{accent}\rule{0.35\textwidth}{0.8pt}}
\end{center}

\vspace{6mm}

% ─── Body ────────────────────────────────────────────────────────────────────
Bonjour,

Je suis étudiant en Réseaux \& Télécommunications à l'UTT et votre offre de stage en sécurité d'infrastructure m'intéresse beaucoup. L'idée de travailler sur une infra auto-hébergée avec Ansible, LXC et Docker, c'est exactement ce que je cherche : pas du cloud abstrait, mais de la vraie admin système avec des machines physiques et des vrais enjeux de sécurité.

Grâce à ma formation et mes projets perso, j'ai des bases solides en administration Linux, Docker, LXC, Bash et les fondamentaux réseau (TCP/IP, DNS, DHCP, MPLS, BGP). J'ai notamment monté des labs OpenVPN avec gestion complète de PKI (CA, génération et distribution de certificats) et construit des environnements de pentesting isolés avec Kali Linux, Nmap et Metasploit. J'ai aussi conçu une infrastructure réseau segmentée (WAN/DMZ/LAN) avec firewall Cisco et VPN lors d'un hackathon à l'UTT. Ce sont des projets où la sécurité n'était pas un bonus mais le sujet principal.

À côté de ça, j'ai fondé novad.app, un SaaS en production avec +1000 utilisateurs (Next.js, PostgreSQL, Vercel). Gérer un service en prod m'a appris à automatiser les déploiements et à comprendre ce qui peut casser quand on touche à l'infra.

Le fait que vous travailliez en IaC avec Ansible et que vous vous tourniez vers Wazuh pour la supervision sécurité, c'est le genre de stack que j'ai envie d'apprendre en conditions réelles. Et le côté logiciel libre me parle : mes projets sont open-source et j'utilise Linux au quotidien.

Je vous incite grandement à consulter mon portfolio où se trouve l'ensemble de mes projets et mon code.

En espérant pouvoir en discuter prochainement.

\vspace{2mm}

Cordialement,

\vspace{3mm}
{\small\color{mutedtext}Amaury DUFRENOT}\\[1.5pt]
{\small\color{mutedtext}Portfolio : \href{https://amaury-dufrenot.me}{amaury-dufrenot.me}}\\[1.5pt]
{\small\color{mutedtext}GitHub : \href{https://github.com/Amorizz}{github.com/Amorizz}}\\[1.5pt]
{\small \href{mailto:amaury.dufrenot@utt.fr}{\color{mutedtext}amaury.dufrenot@utt.fr}\hspace{6pt}{\color{sep}\textbar}\hspace{6pt}\color{mutedtext}+33 6 44 19 68 50}

\end{document}
